\documentclass[../../hippo]{subfiles}
\begin{document}

\subsection{Derivation for Translated Laguerre (HiPPO-LagT)}
\label{sec:derivation-lagt}

We consider measures based on the generalized \emph{Laguerre} polynomials.
For a fixed $\alpha \in \mathbb{R}$, these polynomials $L^{(\alpha)}(t-x)$ are
orthogonal with respect to the measure $x^{\alpha} e^{-x}$ on $[0, \infty)$
(cf.\ \cref{sec:laguerre-properties}).
This derivation will involve tilting the measure governed by another parameter $\beta$.

The result in \cref{thm:legt-lagt} for HiPPO-LagT is for the case $\alpha = 0, \beta = 1$,
corresponding to the basic Laguerre polynomials and no tilting.

\paragraph{Measure and Basis}
We flip and translate the generalized Laguerre weight function and polynomials from $[0, \infty)$ to $(-\infty, t]$.
The normalization is found using equation~\eqref{eq:laguerre}.
\begin{align*}
  \omega(t, x) &=
  \begin{cases}
    (t-x)^\alpha e^{x-t} & \mbox{if } x \le t \\
    0 & \mbox{if } x > t
  \end{cases}
  \quad
  \\
  &= (t-x)^\alpha e^{-(t-x)} \mathbb{I}_{(-\infty, t]}
  \\
  p_n(t, x)
  &= \frac{\Gamma(n+1)^{\frac{1}{2}}}{\Gamma(n+\alpha+1)^{\frac{1}{2}}} L^{(\alpha)}_n(t-x)
\end{align*}


\paragraph{Tilted Measure}


We choose the following tilting $\chi$
\begin{align*}%
  \chi(t, x) &= (t-x)^\alpha \exp\left( -\frac{1-\beta}{2} (t-x) \right) \mathbb{I}_{\Rn{t}} \\
\end{align*}
for some fixed $\beta \in \mathbb{R}$.
The normalization is (constant across all $t$)
\begin{align*}%
  \zeta &= \int \frac{\omega}{\chi^2} = \int (t-x)^{-\alpha} e^{-\beta(t-x)} \mathbb{I}_{\Rn{t}} \dd x
  \\
  &= \Gamma(1-\alpha) \beta^{\alpha-1}
  ,
\end{align*}
so the tilted measure has density
\begin{align*}%
  \zeta(t)^{-1} \frac{\omega^{(t)}}{(\chi^{(t)})^2} &= \Gamma(1-\alpha)^{-1} \beta^{1-\alpha} (t-x)^{-\alpha} \exp\left( -\beta(t-x) \right) \mathbb{I}_{\Rn{t}}
  .
\end{align*}

We choose
\begin{align*}
  \lambda_n = \frac{\Gamma(n+\alpha+1)^{\frac{1}{2}}}{\Gamma(n+1)^{\frac{1}{2}}}
\end{align*}
to be the norm of the generalized Laguerre polynomial $L_n^{(\alpha)}$,
so that $\lambda_n p_n^{(t)} = L_n^{(\alpha)}(t-x)$,
and (following equation~\eqref{eq:framework-g}) the basis for $\nu^{(t)}$ is
\begin{equation}%
  \begin{aligned}
    g_n^{(t)} &= \lambda_n \zeta^{\frac{1}{2}} p_n^{(t)} \chi^{(t)}
    \\
    &= \zeta^{\frac{1}{2}} \chi^{(t)} L_n^{(\alpha)}(t-x)
  \end{aligned}
\end{equation}



\paragraph{Derivatives}

We first calculate the density ratio
\begin{align*}%
  \frac{\omega}{\chi}(t, x) &= \exp \left( -\frac{1+\beta}{2} (t-x) \right) \mathbb{I}_{\Rn{t}}
  .
\end{align*}
and its derivative
\begin{align*}
  \ppt \frac{\omega}{\chi}(t, x)
  &= -\left( \frac{1+\beta}{2} \right)\frac{\omega}{\chi}(t, x) + \exp\left( -\left( \frac{1+\beta}{2} \right)(t-x) \right) \delta_t
  .
\end{align*}

The derivative of Laguerre polynomials can be expressed as linear combinations of
other Laguerre polynomials (cf.\ \cref{sec:laguerre-properties}).
\begin{align*}
  \ppt \lambda_n p_n(t, x)
  &= \ppt L^{(\alpha)}_n(t-x)
  \\
  &= -L^{(\alpha)}_0(t-x) - \dots - L^{(\alpha)}_{n-1}(t-x)
  \\
  &= -\lambda_0 p_0(t, x) - \dots - \lambda_{n-1} p_{n-1}(t, x)
\end{align*}


\paragraph{Coefficient Dynamics}
Plugging these derivatives into equation~\eqref{eq:coefficient-dynamics} (obtained from differentiating the coefficient equation~\eqref{eq:coefficient}), where we suppress the dependence on $x$ for convenience:
\begin{align*}
  \ddt c_n(t)
  &= \zeta^{-\frac{1}{2}} \int f \cdot \left(\ppt \lambda_n p_n^{(t)}\right) \frac{\omega^{(t)}}{\chi^{(t)}}
  \\&\quad + \int f \cdot \left( \zeta^{-\frac{1}{2}} \lambda_n p_n^{(t)} \right) \left( \ppt \frac{\omega^{(t)}}{\chi^{(t)}} \right)
  \\
  &= -\sum_{k=0}^{n-1} \int f \cdot \left( \zeta^{-\frac{1}{2}} \lambda_k p_k^{(t)} \chi^{(t) }\right) \frac{\omega^{(t)}}{(\chi^{(t)})^2}
  \\&\quad - \left( \frac{1+\beta}{2} \right) \int f \cdot \left( \zeta^{-\frac{1}{2}} \lambda_n p_n^{(t)} \right) \frac{\omega^{(t)}}{\chi^{(t)}} + f(t) \cdot \zeta^{-\frac{1}{2}} L_n^{(\alpha)}(0)
  \\
  &= -\sum_{k=0}^{n-1} c_k(t) - \left( \frac{1+\beta}{2} \right) c_n(t) + \Gamma(1-\alpha)^{-\frac{1}{2}} \beta^{\frac{1-\alpha}{2}} \binom{n+\alpha}{n} f(t).
\end{align*}

We then get
\begin{equation}
  \begin{aligned}
    \ddt c(t)
    &= -A c(t) + B f(t)
    \\
    A &=
    \begin{bmatrix}
      \frac{1+\beta}{2} & 0 & \dots & 0 \\
      1 & \frac{1+\beta}{2} & \dots & 0 \\
      \vdots & & \ddots & \\
      1 & 1 & \dots & \frac{1+\beta}{2}
    \end{bmatrix}
    \\
    B &= \zeta^{-\frac{1}{2}} \cdot \begin{bmatrix} \binom{\alpha}{0} \\ \vdots \\ \binom{N-1+\alpha}{N-1} \end{bmatrix}
  \end{aligned}
\end{equation}

\paragraph{Reconstruction}
By equation~\eqref{eq:reconstruction}, at every time $t$, for $x \leq t$,
\begin{align*}
  f(x) \approx g^{(t)}(x)
  &= \sum_{n=0}^{N-1} \lambda_n^{-1} \zeta^{\frac{1}{2}} c_n(t) p_n^{(t)} \chi^{(t)}
  \\
  &= \sum_n \frac{n!}{(n+\alpha)!} \zeta^{\frac{1}{2}} c_n(t) \cdot L^{(\alpha)}_n(t-x) \cdot (t-x)^\alpha e^{\left( \frac{\beta-1}{2} \right)(t-x)}.
\end{align*}

\paragraph{Normalized Dynamics}
Finally, following equations~\eqref{eq:coefficient-dynamics-normalized} and \eqref{eq:reconstruction-normalized} to convert these to dynamics on the orthonormal basis of the normalized (probability) measure $\nu^{(t)}$ leads to the following $\hippo$ operator
\begin{equation}
  \begin{aligned}
    \ddt c(t)
    &= -A c(t) + B f(t)
    \\
    A &=
    -\Lambda^{-1}
    \begin{bmatrix}
      \frac{1+\beta}{2} & 0 & \dots & 0 \\
      1 & \frac{1+\beta}{2} & \dots & 0 \\
      \vdots & & \ddots & \\
      1 & 1 & \dots & \frac{1+\beta}{2}
    \end{bmatrix}
    \Lambda
    \\
    B &= \Gamma(1-\alpha)^{-\frac{1}{2}} \beta^{\frac{1-\alpha}{2}} \cdot \Lambda^{-1}
    \begin{bmatrix} \binom{\alpha}{0} \\ \vdots \\ \binom{N-1+\alpha}{N-1} \end{bmatrix}
    \\
    \Lambda &= \diag_{n \in [N]} \left\{ \frac{\Gamma(n+\alpha+1)^{\frac{1}{2}}}{\Gamma(n+1)^{\frac{1}{2}}} \right\}
  \end{aligned}
\end{equation}
and correspondingly a $\proj_t$ operator:
\begin{equation}
  f(x) \approx g^{(t)}(x)
  = \Gamma(1-\alpha)^{\frac{1}{2}} \beta^{-\frac{1-\alpha}{2}}
  \sum_n c_n(t) \cdot \frac{\Gamma(n+1)^{\frac{1}{2}}}{\Gamma(n+\alpha+1)^{\frac{1}{2}}} \cdot L^{(\alpha)}_n(t-x) \cdot (t-x)^\alpha e^{\left( \frac{\beta-1}{2} \right)(t-x)}.
\end{equation}



\end{document}

